\scanpage{113}

Soit $G$ un schéma \uline{affine} \inserted{\unsure{loc. lib.|loc. lis.}} sur un anneau $A$.

Dans quelle mesure $G$ se récupère-t-il par la connaissance de $({C}_{\deleted{\unsure{f}}}(G),\varphi_G^{\deleted{f}})$, où $C_{\deleted{f}}(G)$ est regardé comme une $\otimes$-catégorie [pour $\otimes$ \ldots], et $\varphi_G^{\deleted{f}}$ comme un $\otimes$-foncteur.

Il faudrait récupérer $G(B)$ pour les $A$-algèbres $B$. Or pour ceci, il suffirait de récupérer $G(A)$. On a un hom. évident

\diagram{\begin{tikzcd}[ampersand replacement=\&,row sep=1.5em]
G(A)\arrow[r] \& \operatorname{Aut}_{\otimes}(\varphi_G^{\deleted{f}})\arrow[d,phantom,"{\deleted{\unsure{\wr|\shortparallel}}}" description]\\
\& \deleted{\operatorname{Aut}_{\otimes}(\varphi_G)}.
\end{tikzcd}}

On présume que c'est injectif : regarder les effets de $g\in G(A)$ \unsure{,} sur \uline{l'algèbre \inserted{affine $\unsure{\varepsilon|\mathscr E|E}$} de} $G$ !

On présume sans doute aussi que c'est surjectif, en regardant encore l'effet d'un $\otimes$-Aut $u$ de $\varphi_G^{\deleted{f}}$ sur $\varepsilon$, \deleted{ce qui donne un \unsure{autom.}} \uline{à détailler}

\uline{Question} À quelle condition une $\otimes$-catégorie \inserted{$A$-linéaire $C_{\deleted{f}}$} munie d'un $\otimes$-foncteur $\varphi^{\deleted{f}}:C_{\deleted{f}}\unsure{-|\to}\operatorname{Mod}_{f}(A)$, \deleted{\unsure{peut|pro}}\unsure{provient-il|provien-il} d'un schéma en groupes affine sur $A$ ???

Il faudrait supposer que $C$ est un Ab-topos, et que $\varphi$ commute aux $\varinjlim$

% archive sheet number: 112
